% s-251-p251-m002; original PDF 251, 252, 253
\Aphorism{s-251-p251-m002}{348}
In Zeiten einer künstlerischen Freizügigkeit ohne Gleichen, einer ungebändigten Sucht zu neuern und zu experimentiren, wo jeder Tag neue Moden, Schlagworte, Raffinements bringt und das mühsam Erquälte von gestern heute schon stereotyp und morgen gar veraltet ist: da hat der productive Künstler nichts mehr zu fürchten als eine allzufeine Sensibilität seines Temperaments und Geschmacks. Es ist gefährlich, heute Geist zu haben: damit geräth man unerbittlich ins allgemeine \Foreign{presto} hinein und fährt am Ende irgendwelchen Abhang hinunter ins Bodenlose. Um nicht versucht zu werden, muss man sich blind stellen; wer die Augen aufthut, ist schon verloren — es giebt zuviel des Verlockenden, Seltsamen, Aufregenden zu schaun. Eine Fülle starker, nicht allzustarker Individualitäten wird heute einfach ausgelöscht und verwischt; sie wissen nicht sie selbst zu bleiben. Zwei, drei Moden mitgemacht, ein wenig naturalisirt, archaisirt, symbolisirt — und mit dem eigenen Wollen und Können ist’s vorbei; fortan reicht die Begabung eben noch zum Schritthalten mit dem jeweilig Neuesten hin, so dass vor lauter Anempfindung und „Con“genialität kein einziges persönliches Wort mehr gehört wird. In Anbetracht, dass es gerade die geistreichen, anschmiegsamen und feinorganisirten Künstler sind, die von der Fluth mit hinweggespült werden, darf man es keinem schaffenden Geiste verargen, wenn er sich gegen die moderne Kunst verschanzt und allein bleibt, auf die Gefahr hin ein Thebaner zu werden: wenn es nur hilft! Denn die Zuckungen der modernen Seele erschüttern selbst den Fernstehenden, und wie man der Muschel nachsagt, dass in ihr das ferne Meer mitbrause, so vibrirt noch im einsamsten Gehäuse des culturfremdesten Einsiedlers etwas von dem, was man den Pulsschlag des modernen Lebens nennt. Es gehört eine freiwillige Beschränktheit und daneben eine Gottesgabe Phlegma dazu, unmodern zu bleiben: aber es ist nöthig, der Instinct grosser Männer beweist es. „Ich muss gestehen, froh zu sein,“ schreibt Keller 1850 an Freiligrath, „dass ich mich durch meine Langsamkeit und Faulheit über diese krankhafte und impotente Periode hinausgerettet habe und zur Vernunft gekommen bin, ohne dergleichen Eseleien zu machen, wozu ich auch grosse Anlagen hatte.“ Das war Krankhaftigkeit von 1850, noch lange nicht Krankhaftigkeit von 1895; damals spukte ein „pikanter Byron’scher Atheismus“ und Heine’scher Selbstwiderspruch, aber was spukt nicht alles heute? Der Symbolismus. Der Neo-Idealismus. Der Spiritismus. Der Japanismus. Der Satanismus. Das Geschlechtliche, nicht mehr in vierzig, sondern in vierhundert Abarten. Die Mystik, nicht mehr katholischer Kirchen, sondern der Riesentempel von Elephante. Die Neurasthenie, die von gelber chinesischer Seide hallucinirt, während „in rüttelnden Wogen Krämpfe über ihre Knochen rieselten.“ Die Sprache, nicht mehr gallisirend, sondern zugleich keltisirend, zigeunernd, in allen europäischen Nüancen flimmernd. Die Farbe: schwindend, unfassbar, übersinnlich — oder schreiend conträr: grüner Himmel auf violetten Wiesen. Die Musik: nur noch Chromatik, Stimmungsgewühl, lechzend, stammelnd, rasend, Brunst und keuchender Athem. In diesem Wirbel leben und taumeln wir Modernen. Zwei, drei Grosse stehen abseits: es sind die Keller-Naturen, Schweizer von Temperament, vielleicht von Abstammung. Der Rest ruinirt sich \Dots{3}
\par

